\section{Permissioned Blockchains and \ac{PoA} Consensus Mechanisms}

Based on the participation and access control blockchain network can be classified in to two categories.  \cite{und2026commledger}. They are permissionless architecture (public) and permissioned architecture. Bitcoin and Early Ethereum networks are permissionless blockchains. They allowed unrestricted participation and relyed on consensus mechanism like \ac{PoW} or Proof of Stake (\ac{PoS})\cite{mdpi2023survey}. These architectures and consensus mechanisms emphasized decentralization among participants. But they often suffer from the high latency of the network, limited throughput and unpredictable transaction costs due to congestion. These problems make the traditional permissionless PoW or PoS consensus mechanisms unsuitable for governance and crowdsourced applications\cite{arxiv2024performance}.

While permissionless blockchains are unsuitable for governance related applications, permissioned blockchains restrict the network participation to authorized entities. This enables more efficient consensus mechanisms and greater control over governance data \cite{und2026commledger}. This known identities of participants allow permissioned systems to achieve higher performance. And support the regulatory compliances, data privacy and structured governance models as well cite{mdpi2023unleashing}.

Other than \ac{PoW} and \ac{PoS} consensus mechanisms, \ac{PoA} offer significant advantages over permissioned blockchains. \ac{PoW} requires energy intensive computation. \ac{PoS} relies on the economic stake and introduce wealth centralization risks. But on the other hand \ac{PoA} depends on the identity and the reputation of pre approved validators\cite{researchgate2020survey}, \cite{zebpay2026poa}. Therefore this \ac{PoA} consensus mechanism enables high transaction throughput and low latency. And also it mitigate the issues such as the "nothing at stake" problem by off chain legal and reputational accountability \cite{echo2026poa}, \cite{grokipedia2026poa}.

The government systems should have accountability, performance predictability and data confidentiality. The \ac{PoA} is well suited consensus mechanism for that \cite{onekey2026poa}. By using \ac{PoA} we can have the ability to identify and regulate validator entities to align blockchain operations. It's because we have the existing legal frameworks that ensure institutional responsibility. Hence, can avoid malicious behavior \cite{hep2023governance}. In Addition to that PoA avoids the indeterministic transaction fees and support controlled validation environments. This make them viable for handling sensitive governmental data under regulatory constraints\cite{zebpay2026poa},\cite{mdpi2023unleashing}. As a result, \ac{PoA} emerges as a practical consensus mechanism for permissioned blockchain-based reporting systems in local governance contexts.

